%----------------------------------------------------------------------------------------
%	PACKAGES AND OTHER DOCUMENT CONFIGURATIONS
%----------------------------------------------------------------------------------------

\documentclass[12pt]{article}
\usepackage[danish]{babel}
\usepackage[numbered,final]{mcode}
\usepackage[utf8x]{inputenc}
\usepackage{amsmath}
\usepackage{graphicx}
\usepackage[colorinlistoftodos]{todonotes}
\usepackage[toc,page]{appendix}
%\usepackage{float}
\usepackage{floatrow} % used for adding "Source" to pictures
\usepackage{hyperref} % used for hyperlinks
\usepackage[all]{hypcap}
\usepackage{bm} % used for bold inline matj
\usepackage{lipsum} % used for lorem lipsum
\usepackage[final]{pdfpages} % used for including PDF's
\usepackage{geometry}
\usepackage{listingsutf8}
\usepackage{listings}
\usepackage{color} %red, green, blue, yellow, cyan, magenta, black, white
\definecolor{mygreen}{RGB}{28,172,0} % color values Red, Green, Blue
\definecolor{mylilas}{RGB}{170,55,241}

\hypersetup{colorlinks=true, linkcolor=black}

% Page margins
\geometry{verbose,tmargin=1in,bmargin=1in,lmargin=1in,rmargin=1in,headsep=0.35in}


\begin{document}

\begin{titlepage}



\newcommand{\HRule}{\rule{\linewidth}{0.5mm}} % Defines a new command for the horizontal lines, change thickness here
\setlength{\topmargin}{0in}
\centering % Center everything on the page

%----------------------------------------------------------------------------------------
%	HEADING SECTIONS
%----------------------------------------------------------------------------------------
\textsc{\LARGE Aarhus universitet}\\[1.5cm] % Name of your university/college
\textsc{\Large Introduktion til digital signalbehandling}\\[0.5cm] % Major heading such as course name
\textsc{\large 3. Semester}\\[0.5cm] % Minor heading such as course title

%----------------------------------------------------------------------------------------
%	TITLE SECTION
%----------------------------------------------------------------------------------------

\HRule \\[0.4cm]
{ \huge \bfseries DSA Case 2}\\ % Title of your document
\HRule \\[1cm]
 
%----------------------------------------------------------------------------------------
%	AUTHOR SECTION
%----------------------------------------------------------------------------------------

\begin{minipage}{0.4\textwidth}
	\begin{flushleft} \large
		\emph{Gruppemedlemmer:}\\
		Gustav A. Gammelgaard \\
		Stinus Lykke Skovgaard \\
		Tim Hede Stenholt \\ [0.5cm]
	\end{flushleft}
\end{minipage}
~
\begin{minipage}{0.4\textwidth}
	\begin{flushright} \large
		\emph{AUid:} \\
		au538293\\
		au520659\\
		au543518\\ [0.5cm]
	\end{flushright}
\end{minipage}\\[5cm]

%----------------------------------------------------------------------------------------
%	LOGO SECTION
%----------------------------------------------------------------------------------------

\includegraphics[scale=0.5]{Img/logo.jpg}\\[1cm]

%----------------------------------------------------------------------------------------
%	DATE SECTION
%----------------------------------------------------------------------------------------

{\large \today}\\[0.5cm] % Date, change the \today to a set date if you want to be precise


\vfill % Fill the rest of the page with whitespace

\end{titlepage}

\newpage
\tableofcontents
\newpage
\listoffigures
\newpage

\hypersetup{linkcolor=blue}

\section{Problem beskrivelse}
\begin{flushleft}
	
	
	Denne case omhandler måden at forbedre et støjfyldt signal med et midlingsfilter. 
	
	
	Denne case omhandler filtrering af et lydklip. Dette lydklip består af en mand der taler mens en tone spiller inderover. Denne tone skal filtreres fra, så det kun er er mandens stemme der er tilbage. 
	Dette er implementeret på et Blackfin development kit. Filteret er et notch-filter, hvilket er et meget smalt filter.
\end{flushleft}

\section{Opgave 1}

\section{Opgave 2}

\section{Opgave 3}

\section{Matlab kode}	

\end{document}